% =============================================================================
%  PARTIE II — Comprendre Icare
% =============================================================================

\chapter{Le DMX512 en bref}
\label{ch:dmx}

Cette partie explique \emph{pourquoi} le code fait ce qu'il fait. On y revoit le
protocole DMX, la trame de l'adaptateur, la circulation des données et les
algorithmes des effets.

\section{Un univers de 512 canaux}

DMX512 est un protocole de scène très simple. Un \emph{univers} contient
\textbf{512 canaux}, chacun sur un octet, donc une valeur de \textbf{0 à 255}.
Onze bits sont envoyés par canal : une impulsion de \emph{start bit}, huit bits de
données, deux \emph{stop bits}. Le tout circule sur un bus RS-485 à
\textbf{250 000 bauds}.

Deux propriétés importantes :

\begin{itemize}[leftmargin=1.4em]
  \item \textbf{Aucun retour.} Le contrôleur parle, les projecteurs écoutent.
        On ne sait jamais si un projecteur a compris autre chose que ce qu'on
        voulait.
  \item \textbf{Image continue.} La même \emph{image} (les 512 valeurs) est
        renvoyée en boucle. Un projecteur qui cesse de recevoir se met en
        sécurité (souvent après quelques dixièmes de seconde).
\end{itemize}

Les valeurs n'ont pas de sens universel : c'est le projecteur qui décide, selon
son \emph{profil}. Un PAR RGB courant utilise trois canaux (\emph{intensité},
\emph{rouge}, \emph{vert}, \emph{bleu} sur certains modèles, parfois plus).

\section{Le plan de canaux d'\Icare{}}

Pour rester simple, \Icare{} pilote un unique projecteur RGB selon un plan de
canaux fixe, hérité de l'ancien prototype. Ces constantes vivent dans
\file{webapp/config.py}, classe \code{Channels}.

\begin{table}[H]
  \centering
  \begin{tabular}{@{}clp{8cm}@{}}
    \toprule
    \textbf{Canal} & \textbf{Constante} & \textbf{Rôle} \\
    \midrule
    0 & \code{MASTER} & Toujours à 255 pour «{}\/ouvrir\/»{} la sortie. \\
    1 & \code{AMPLITUDE} & Luminosité globale (0--255). \\
    3 & \code{PAN} & Orientation horizontale. \\
    5 & \code{TILT} & Orientation verticale. \\
    8 & \code{SHUTTER} & Toujours à 0 (pas d'effet de shutter). \\
    9 & \code{BASE\_COLOR\_DEFAULT} & R du projecteur, G en 10, B en 11. \\
    \bottomrule
  \end{tabular}
  \caption{Plan de canaux utilisé par le contrôleur.}
  \label{tab:channels}
\end{table}

Le \emph{canal de base} est modifiable : les canaux R, G, B sont alors
\code{base}, \code{base+1} et \code{base+2}. \Icare{} accepte une base de 0 à
509, afin qu'il reste toujours trois canaux disponibles.

\chapter{L'adaptateur USB-DMX}
\label{ch:enttec}

Le PC n'a pas de port DMX natif : on passe par un boîtier USB dont l'un des
plus répandus est l'\emph{Enttec USB Pro}. Le dialogue se fait à 250 000 bauds,
en 8 bits, sans parité, 2 bits de stop (8N2). \Icare{} occupe la classe
\code{DMX\_Icare} de \file{pydmxIcare.py}.

\section{La trame Enttec USB Pro}

Chaque envoi est encapsulé dans une petite trame :

\begin{center}
\begin{tikzpicture}[font=\ttfamily\small,every node/.style={draw,minimum height=0.9cm,minimum width=1.5cm,align=center}]
  \node[fill=icarenote] (a) {0x7E};
  \node[right=0pt of a,fill=icareblue!12] (b) {0x06};
  \node[right=0pt of b,fill=icareblue!12] (c) {LSB};
  \node[right=0pt of c,fill=icareblue!12] (d) {MSB};
  \node[right=0pt of d,fill=icarenote] (e) {0x00};
  \node[right=0pt of e,fill=green!12,minimum width=2.6cm] (f) {512 octets};
  \node[right=0pt of f,fill=icarenote] (g) {0xE7};
\end{tikzpicture}
\end{center}
\legende{Trame Enttec USB Pro : début, étiquette «{}Output Only Send DMX Packet
Request{}», longueur de charge utile en petit-boutiste, code de départ DMX,
données, fin.}

Détail des champs :

\begin{itemize}[leftmargin=1.4em]
  \item \code{0x7E} : octet de début de trame (\emph{START\_BYTE}).
  \item \code{0x06} : étiquette «{}envoi seul de paquet DMX{}»
        (\emph{LABEL\_DMX\_TX}).
  \item \textbf{LSB, MSB} : longueur de la charge utile = 513 octets
        (1 code de départ + 512 canaux), soit \code{0x01 0x02}.
  \item \code{0x00} : code de départ DMX standard.
  \item \textbf{512 octets} : les valeurs des canaux.
  \item \code{0xE7} : octet de fin de trame (\emph{END\_BYTE}).
\end{itemize}

Dans le code, l'en-tête est précalculé une fois pour toutes sous forme de
tableau NumPy, puis concaténé aux 512 canaux à chaque envoi :

\begin{lstlisting}[language=Python,caption={Construction et envoi d'une trame dans \file{pydmxIcare.py}.}]
payload_len = NUM_CHANNELS + 1  # 1 start code + 512 channels
self.start_byte = np.array([
    START_BYTE, LABEL_DMX_TX,
    payload_len & 0xFF, (payload_len >> 8) & 0xFF,
    0x00,  # DMX start code
], np.uint8)
...
frame = np.concatenate((self.start_byte, self.data, self.end_byte)).tobytes()
self.com.write(frame)
\end{lstlisting}

\section{Détection du port}

Si aucun port n'est fourni, \Icare{} parcourt les ports série et retient le
premier dont la description contient \og USB \fg{} :

\begin{lstlisting}[language=Python,caption={Auto-détection du port série.}]
for port_no, description, _address in serial.tools.list_ports.comports():
    if "USB" in description:
        return str(port_no)
return None
\end{lstlisting}

En cas d'échec d'écriture, le port est refermé proprement plutôt que de laisser
la boucle échouer en rafale.

\chapter{La chaîne logicielle, de bout en bout}
\label{ch:chaine}

Voici comment un mouvement de curseur finit en photons.

\begin{figure}[H]
  \centering
  \begin{tikzpicture}[
      node distance=8mm,
      font=\small,
      every node/.style={draw=icodeln,fill=icodebg,rounded corners,align=center,inner sep=6pt},
      flow/.style={-Latex,draw=icareblue,thick},
    ]
    \node (browser) {Navigateur\\(curseur, effets, audio)};
    \node[right=1.5cm of browser] (flask) {Flask\\API REST};
    \node[right=1.5cm of flask] (dmx) {\code{DMXController}\\thread de sortie};
    \node[below=1.1cm of dmx] (serial) {\code{pydmxIcare}\\port série};
    \node[left=1.5cm of serial] (adapter) {Adaptateur\\USB-DMX};
    \node[left=1.5cm of adapter] (fixtures) {Projecteurs};
    \draw[flow] (browser) -- node[above,draw=none,fill=none,font=\scriptsize]{HTTP JSON} (flask);
    \draw[flow] (flask) -- node[above,draw=none,fill=none,font=\scriptsize]{appel Python} (dmx);
    \draw[flow] (dmx) -- node[right,draw=none,fill=none,font=\scriptsize]{trame} (serial);
    \draw[flow] (serial) -- node[above,draw=none,fill=none,font=\scriptsize]{USB} (adapter);
    \draw[flow] (adapter) -- node[above,draw=none,fill=none,font=\scriptsize]{DMX512} (fixtures);
    \draw[flow,dashed] (fixtures.south) to[out=-90,in=-90,distance=3.2cm] node[below,draw=none,fill=none,font=\scriptsize]{lumières (aucun retour)} (browser.south);
  \end{tikzpicture}
  \caption{Chaîne complète, du navigateur aux projecteurs.}
  \label{fig:chain}
\end{figure}

Les étapes :

\begin{enumerate}[leftmargin=1.6em]
  \item Le navigateur (\file{app.js}) envoie par exemple
        \api{POST /api/experiment/set-values} avec \code{\{r, g, b, amp, pan, tilt\}}.
  \item La route correspondante (\file{webapp/api\_experiment.py}) appelle
        \code{dmx.set\_values(...)}.
  \item \code{DMXController} borne les valeurs (0--255) et, selon le mode actif,
        écrit la \emph{frame} courante.
  \item La boucle de fond (thread \code{dmx-out}) rafraîchit la sortie en
        continu (environ 20 Hz en mode stable).
  \item \code{pydmxIcare} assemble la trame Enttec et l'écrit sur le port série.
  \item L'adaptateur traduit cela en signal DMX512 sur la ligne.
\end{enumerate}

\section{Pourquoi un thread ?}

L'écriture série et les effets sont des tâches \emph{continues} qui ne doivent
pas bloquer les requêtes HTTP. \Icare{} isole donc la sortie DMX dans un thread
démon, et le lecteur de séquence dans un second thread. Les routes HTTP ne font
que modifier l'état partagé (couleurs, paramètres de mode) ; le thread s'occupe
de le refléter sur le fil.

\section{L'état partagé}

Le contrôleur maintient :

\begin{itemize}[leftmargin=1.6em,itemsep=2pt,topsep=3pt]
  \item \code{r, g, b, amp, pan, tilt} : les valeurs de sortie ;
  \item \code{enabled, strobe\_enabled, random\_enabled, effect\_enabled,
        audio\_enabled} : un seul mode actif à la fois ;
  \item \code{overlay} : le strobe superposé, indépendant du mode ;
  \item \code{random\_params, effect\_params, base\_color\_channel}.
\end{itemize}

\chapter{Les effets procéduraux}
\label{ch:effets}

Quatre effets sont générés côté serveur, dans la boucle de sortie. Le temps
écoulé depuis le démarrage, multiplié par la \emph{vitesse}, sert de paramètre
d'animation.

\begin{table}[H]
  \centering
  \begin{tabular}{@{}lp{10.5cm}@{}}
    \toprule
    \textbf{Effet} & \textbf{Principe} \\
    \midrule
    \code{rainbow} & Teinte qui tourne : $h = (60\,t) \bmod 360$ degrés, saturation
      et valeur maximales (conversion HSV$\rightarrow$RGB). \\
    \addlinespace
    \code{fire} & Orange/rouge vacillant : teinte autour de 18\textdegree, valeur
      aléatoire entre 0,55 et 1,0 pour l'effet de flamme. \\
    \addlinespace
    \code{police} & Alternance rouge / bleu à environ 4 Hz
      ($\lfloor 4t \rfloor \bmod 2$). \\
    \addlinespace
    \code{chase} & Chenillard qui parcourt R, G, B à environ 3 pas par seconde. \\
    \bottomrule
  \end{tabular}
  \caption{Effets procéduraux intégrés.}
\end{table}

La conversion HSV vers RGB est implémentée à la main, sans dépendance externe :

\begin{lstlisting}[language=Python,caption={Le cœur de \code{\_hsv\_to\_rgb}.}]
c = v * s
x = c * (1 - abs(((h / 60) % 2) - 1))
m = v - c
if   h <  60: r1, g1, b1 = c, x, 0
elif h < 120: r1, g1, b1 = x, c, 0
...
return int((r1 + m) * 255), int((g1 + m) * 255), int((b1 + m) * 255)
\end{lstlisting}

\section{Le strobe en superposition}

Le strobe d'overlay est à part : il ne remplace pas le mode courant, il passe
par-dessus. À chaque image, on évalue une phase
$\phi = (t \cdot f) \bmod 1$ et on éteint la sortie quand $\phi$ dépasse le
rapport cyclique :

\begin{equation*}
\text{sortie} =
\begin{cases}
\text{couleur courante} & \text{si } \phi < \text{duty}, \\
\text{noir (ou 0 sur le canal choisi)} & \text{sinon.}
\end{cases}
\end{equation*}

Si un canal DMX est précisé, seul ce canal clignote ; sinon toute l'image
alterne. L'overlay s'arrête tout seul à la fin de la durée demandée.

\chapter{Le mode audio-réactif}
\label{ch:audio}

Ce mode transforme le microphone du navigateur en capteur de rythme. Toute
l'analyse se fait côté client ; le serveur ne reçoit que le résultat.

\begin{figure}[H]
  \centering
  \begin{tikzpicture}[
      node distance=9mm,
      font=\scriptsize,
      every node/.style={draw=icodeln,fill=icodebg,rounded corners,align=center,inner sep=5pt},
      flow/.style={-Latex,draw=icareblue,thick},
    ]
    \node (mic) {Micro\\\code{getUserMedia}};
    \node[right=0.35cm of mic] (an) {\code{Analyser}\\\code{fftSize=1024}};
    \node[right=0.35cm of an] (rms) {RMS\\\code{getByteTimeDomainData}};
    \node[right=0.35cm of rms] (lvl) {niveau $\times$ gain\\borne à $[0,1]$};
    \node[right=0.35cm of lvl] (api) {\api{POST}\\\code{/audio/push}};
    \node[below=1.2cm of api] (dmx) {\code{amp = niveau $\times$ 255}};
    \draw[flow] (mic) -- (an);
    \draw[flow] (an) -- (rms);
    \draw[flow] (rms) -- (lvl);
    \draw[flow] (lvl) -- (api);
    \draw[flow] (api) -- (dmx);
  \end{tikzpicture}
  \caption{Chaîne audio-réactive.}
  \label{fig:audio}
\end{figure}

\paragraph{Du son au niveau.} On échantillonne la forme d'onde, on recentre les
octets autour de 128, puis on calcule la valeur quadratique moyenne (RMS) :

\begin{equation*}
\mathrm{RMS} = \sqrt{\frac{1}{N}\sum_{i=0}^{N-1} \left(\frac{d_i - 128}{128}\right)^2 },
\qquad
\text{niveau} = \min\!\big(1,\ \mathrm{RMS} \times \text{gain}\big).
\end{equation*}

Le curseur \emph{Sensibilité} règle le \code{gain} entre 0,5 et 6.

\paragraph{Fréquence d'envoi.} La boucle d'animation tourne à la fréquence de
l'écran, mais l'envoi réseau est limité à environ 30 Hz pour ne pas saturer le
serveur. La couleur choisie accompagne le niveau.

\paragraph{Côté serveur.} \code{audio\_push(niveau, couleur)} met à jour
l'amplitude à \code{niveau $\times$ 255}, applique éventuellement la couleur, et
renvoie l'image. Le mode ne génère rien tout seul : sans échantillons, l'état
ne change pas.
